% ============================================================
% The Cross-Alignment Matrix
% ============================================================

\documentclass[12pt]{amsart}

\usepackage{amsmath,amssymb,amsthm}
\usepackage{booktabs}
\usepackage{microtype}
\usepackage{url}

\newtheorem{theorem}{Theorem}
\newtheorem{lemma}[theorem]{Lemma}
\newtheorem{proposition}[theorem]{Proposition}
\newtheorem{corollary}[theorem]{Corollary}
\newtheorem{conjecture}[theorem]{Conjecture}
\theoremstyle{definition}
\newtheorem{definition}[theorem]{Definition}
\theoremstyle{remark}
\newtheorem{remark}[theorem]{Remark}
\newtheorem{observation}[theorem]{Observation}

\title{The Cross-Alignment Matrix}

\author{Alexander S.\ Petty}

\email{alexander.petty@gmail.com}
\date{April 2021}
\subjclass[2020]{11A63, 11B83, 42A16}

\begin{document}
\begin{abstract}
We introduce the cross-alignment matrix of a fractional
field: the symmetric matrix $\mathbf{A}(n)$ whose
$(i,j)$ entry is the position-wise digit-match
proportion between $i/n$ and $j/n$. For
digit-partitioning primes, $\mathbf{A}(p) = I$.
For $n = ps$ with $p$ digit-partitioning and $s$
smooth, $\mathbf{A}$ decomposes into $p$ blocks of
rank~$1$, giving $\operatorname{rank}(\mathbf{A}) = p$.

The pairwise alignment $\sigma(n)$ is the mean
off-diagonal entry. For primitive-root primes, the
matrix is permutation-similar to a circulant whose
eigenvalues are the DFT of the cyclic autocorrelation.
\end{abstract}
\maketitle

% ============================================================
\section{Introduction}

Scalars measure one thing at a time. The matrix
measures everything at once.

The preceding papers measured the fractional field
$\mathcal{F}(n) = \{k/n : 1 \le k \le n{-}1\}$ through
$\alpha(n)$, $\sigma(n)$, and $F(n)$. This paper
introduces the object they are projections of: the
cross-alignment matrix.

% ============================================================
\section{Definition}

\begin{definition}
The \emph{cross-alignment matrix} of $n$ in base $b$ is the
$(n{-}1) \times (n{-}1)$ symmetric matrix $\mathbf{A}(n)$ with
entries
\[
A_{ij} = a(i/n,\, j/n),
\]
where $a(i/n,\, j/n)$ is the position-wise digit-match proportion
between the repetends of $i/n$ and $j/n$. Terminating pairs
contribute~$1$; pairs with different cycle lengths contribute~$0$;
the diagonal satisfies $A_{ii} = 1$.
\end{definition}

The matrix is symmetric by construction ($A_{ij} = A_{ji}$)
and positive semidefinite, since within each period class
the submatrix is a scalar multiple of an all-ones matrix
(which is PSD), and the direct sum of PSD matrices is
PSD~\cite{horn}.

% ============================================================
\section{Recovery of Scalar Invariants}

The pairwise alignment $\sigma$ is recovered directly
from $\mathbf{A}$:
\[
\sigma(n) = \binom{n-1}{2}^{-1} \sum_{i<j} A_{ij}
\qquad \text{(mean off-diagonal entry)}.
\]

The alignment $\alpha(n)$ is not determined by
$\mathbf{A}$ alone, because $\alpha$ counts terminating
fractions as aligned with $1/n$ (contributing $1$),
while $\mathbf{A}$ assigns $0$ to pairs with different
cycle lengths. The matrix captures the non-terminating
pairwise structure. The alignment $\alpha$ additionally
requires the terminating-fraction convention from the
earlier papers.

% ============================================================
\section{Spectral Structure for Digit-Partitioning Primes}

\begin{theorem}\label{thm:dp-identity}
If $p$ is a digit-partitioning prime in base $b$ (i.e., $p \le b{+}1$),
then $\mathbf{A}(p) = \mathbf{I}_{p-1}$, the identity matrix.
\end{theorem}

\begin{proof}
By the digit-partitioning property, no two distinct fractions
$i/p$ and $j/p$ share a digit at any common position. Thus
$A_{ij} = 0$ for $i \ne j$ and $A_{ii} = 1$.
\end{proof}

All eigenvalues equal~$1$. The field is maximally non-degenerate:
every fraction is structurally independent from every other.

% ============================================================
\section{Spectral Structure for Non-DP Primes}

\begin{observation}\label{obs:non-dp-spectrum}
For a prime $p$ not dividing $b$ with $p > b{+}1$,
the eigenvalue spectrum of $\mathbf{A}(p)$ reflects
the coset structure of $\langle b \rangle$ in
$(\mathbb{Z}/p\mathbb{Z})^{*}$.

When $b$ is a primitive root modulo $p$, the matrix
is permutation-similar to a circulant~\cite{davis}
(reordering rows and columns by discrete logarithm),
with eigenvalues given by the DFT of the cyclic
autocorrelation.

When there are exactly two cosets ($L = (p{-}1)/2$),
computation shows the spectrum has two eigenvalue
levels whose values depend on the cross-coset
digit-match count.
\end{observation}

In base~$10$: $p = 13$ ($L = 6$, two cosets) has eigenvalues
$4/3$ and $2/3$, each with multiplicity~$6$. The prime $p = 17$
($L = 16$, primitive root) has nine distinct eigenvalues ranging
from $0.25$ to $1.75$.

% ============================================================
\section{Block Structure}

\begin{proposition}\label{prop:block}
For $n = ps$ with $p$ a digit-partitioning prime and
$s$ a $b$-smooth integer, the cross-alignment matrix is
permutation-similar to
\[
\mathbf{A}(ps) \cong J_{s-1} \oplus
\underbrace{J_s \oplus \cdots \oplus J_s}_{p-1},
\]
where $J_m$ is the $m \times m$ all-ones matrix.
In particular, $\operatorname{rank}(\mathbf{A}(ps)) = p$.
\end{proposition}

\begin{proof}
The $s - 1$ terminating fractions form one block
(pairwise alignment $1$). The $(p{-}1)$ exact
repetend classes, each of size $s$, form the remaining
blocks (within-class alignment $1$, cross-class $0$
by the digit-partitioning property). Each all-ones
block $J_m$ has rank $1$.
\end{proof}

For digit-partitioning $n = ps$, the rank is always $p$,
regardless of $s$. The null space grows with the smooth
factor because each repetend class of size $s$ contributes
rank $1$ (all members are structurally identical in
alignment space).


% ============================================================
\section{Remarks}

\subsection*{What the matrix captures}

The pairwise alignment $\sigma$ is directly recoverable
from $\mathbf{A}$. The alignment $\alpha$ requires the
additional terminating-fraction convention. For
digit-partitioning primes, $\mathbf{A}(p) = I$
characterizes the property completely.


\subsection*{Open directions}

Whether the eigenvalue spectrum of $\mathbf{A}(n)$ admits a
closed-form description in terms of the prime factorization of $n$
and the digit function is open. The connection
to Dirichlet characters~\cite{iwaniec}, suggested by the coset-level
structure of the spectrum, is pursued in the papers that follow.

% ============================================================
\begin{thebibliography}{9}

\bibitem{davis}
P.~J.~Davis,
\emph{Circulant Matrices},
2nd ed., Chelsea, 1994.

\bibitem{horn}
R.~A.~Horn and C.~R.~Johnson,
\emph{Matrix Analysis},
2nd ed., Cambridge University Press, 2013.

\bibitem{iwaniec}
H.~Iwaniec and E.~Kowalski,
\emph{Analytic Number Theory},
Amer.\ Math.\ Soc., 2004.

\end{thebibliography}

\end{document}
